\chapter{Conclusion}\label{chap:conclusion}

This thesis presents a novel system design architecture for interactive musical prediction, which combines the strengths of both machine learning and rule-based approaches to generate coherent and expressive musical responses in real-time. By integrating Monte Carlo Tree Search (MCTS) with a Mixture Density Recurrent Neural Network (MDRNN), we have developed a system which combines the strengths of heuristic based music generation in structural and musical cohensive with the strengths of neural netowrk based systems in creative and expressive range, while mitigating the inverse weaknesses of each approach when used individually. 

TODO edit:
This thesis successfully presents a novel interactive musical prediction system that significantly enhances the capabilities of the existing IMPSY framework by integrating Monte Carlo Tree Search (MCTS) with a Mixture Density Recurrent Neural Network (MDRNN). Our primary objective was to improve the coherence and expressiveness of machine-generated musical responses in a real-time call-and-response setting, moving beyond the limitations of purely random sampling from a Gaussian Mixture Model.

\subsection{Summary of Key Findings}

TODO (maybe make part before that a summary of what IMPSY IS and how it was created, and this about the findings)
TODO edit whole subsection

The evaluation, conducted using both quantitative and qualitative metrics across different datasets, unequivocally demonstrates the effectiveness of our MCTS-enhanced approach. We observed significant improvements in prediction accuracy and, crucially, in the qualitative aspects of playability and expressiveness, especially within structured improvisational contexts. While the benefits were less pronounced in freeform improvisations, the system still produced more coherent outputs compared to the traditional MDRNN method, without sacrificing creativity.

This research serves as a robust proof of concept for the integration of search algorithms and heuristic-based guidance within machine learning prediction systems for musical applications. The customizability and generalisability of our heuristic components mean that the system can be adapted to various musical styles, instruments, and performance goals, offering a flexible foundation for future developments in interactive artistic systems. The implications extend beyond music, suggesting a pathway for more meaningful and controllable human-AI collaboration in other creative domains. Our work contributes to a deeper, more accessible, and human-centered approach to digital music interaction, addressing the challenging balance between structured musical coherence and the flexibility of creative improvisation.

Moving forward, future work could explore expanding the system's ability to identify and utilize repeated musical patterns and motifs, which are crucial for improvisational coherence. Further human interaction elements, such as a user interface for real-time parameter modification and the ability to train new performance parameters without coding, would also greatly enhance the system's usability and creative potential. This thesis lays a strong foundation for future research in fostering more intelligent and musically sensitive AI collaborators.

\subsection{Answers to the Research Questions}

In the introduction of this thesis, three research questions were posed, which we will now summarise the answers to, which is longer form are distributed throughout the thesis.

\textbf{1. What are the strengths and weaknesses of heuristic and machine learning based approaches in NIMEs, and how can they be effectively combined?}

TODO

\textbf{2. How can an effective heuristic be designed for MCTS-IMPSY to intelligently guide the search process?}

TODO

\textbf{6. Can MCTS-IMPSY improve on the prediction accuracy and playability of MDRNN-IMPSY, while maintaining creativity?}

TODO




